\scanpage{66}

% top, lines 1–3
et la condition que les $u_n$ définissent bien un hom de complexes est la ``condition de Stokes''

% upper, Stokes relation beside its square
\diagram{$\unsure{K_n|k_n|\kappa_n}\partial_{n+1}=K_{n+1}d:\quad$
\begin{tikzcd}[ampersand replacement=\&]
A^n_{n+1}\arrow[r,"\partial_{n+1}"]\arrow[d,"d^n"'] \& A^n_n\arrow[d,"K_n"]\\
A^{n+1}_{n+1}\arrow[r,"K_{n+1}"'] \& M
\end{tikzcd}}

% upper middle, below square
Sous ces conditions, pour tt $X_*$\inserted{$\in\operatorname{Ob}\widehat{\Delta}$}, on a un hom. \deleted{\unsure{ca}} déduit de $u$

% middle, homomorphism and vertically identified target
\diagram{\begin{tikzcd}[ampersand replacement=\&,column sep=small,row sep=tiny]
\Gamma(X_*,A^*_{X_*})\arrow[r] \& C^{\unsure{*|\bullet}}(X_*,M_*)\\[-1ex]
\& \arrow[u,equal]\deleted{\unsure{C}}\Gamma(X_*,C^{*}_{\deleted{\unsure{X}}*M\deleted{\unsure{*}}}\mid X_\bullet)
\end{tikzcd}}

% middle, derived category
qui représente l'hom dans la cat. dérivée

% middle lower, derived-category map
\[
\Gamma(X_*,A^*_{X_*})\mathrel{\unsure{\relbar|\longrightarrow}}\mathbb{R}\Gamma_{X_\bullet}(M)
\]

% lower, before second square
associé à $u_*$. On aura alors commutativité dans

% lower, canonical square with cancelled first letter
\diagram{\begin{tikzcd}[ampersand replacement=\&,column sep=small]
\deleted{\mathbb{R}}\Gamma(X_\bullet,A^*_{X_*})\arrow[r]\arrow[d] \& C^*(X_\bullet,M)=\Gamma(X_\bullet,C^*_{M}\mid\unsure{X_\bullet|X_*})\arrow[d,"\sim"]\\
\mathbb{R}\Gamma(X_\bullet,A^*_{X_\bullet})\arrow[r] \& \mathbb{R}\Gamma(X_*,M)
\end{tikzcd}}

% foot, conditional sentence continues
Si $A^{*}$ est une résolution de $M_*$, et $K_0:A_0^0\longrightarrow M$ induit l'identité sur $M\subset A^0_0$, alors la \unsure{dernière|deuxième} flèche horizontale est aussi un iso, et

% lower left
% archive sheet number: 65
